\SourceLeaf{8}
and the distinguished \textsc{Legendre} succeeded in completing this most difficult part of the problem (\textit{Recherches sur l’attraction des sphéroïdes homogènes, Mémoires présentés à l’acad. roi. des sc. T.} X.)

\textsc{Mac Laurin} had already begun the most general investigation of the attraction of spheroids that are not generated by revolution, but whose sections by any plane are ellipses. He had stopped, however, at the attraction on points lying on one of the three axes. In the paper just cited, the distinguished \textsc{Legendre} had already conjectured by induction the principal theorem on which the most general solution of the problem chiefly rests. But the distinguished \textsc{Laplace} was the first to succeed in proving everything rigorously and so making the solution complete in every respect (\textit{Hist. de l’acad. roi. des sc. de Paris} 1782; the same solution is repeated in \textit{Théorie du mouvement et de la figure elliptique des planètes} and \textit{Mécanique céleste} Vol. 2).

\expandafter\def\csname HistoricalNoteText008\endcsname{The distinguished \textsc{Lagrange}, for example, judges these two solutions as follows. \textit{Their solutions can only be regarded as masterpieces of analysis. Yet one may still wish for a more direct and simpler solution, and the continuing advances of analysis give reason to hope for one.} Nouv. Mém. de l’acad. de Berlin 1793. p. 263.}
Everyone will surely admire the outstanding elegance and subtle ingenuity of this solution by the distinguished \textsc{Laplace}. Yet the very subtlety and admirable skill with which he overcame the formidable difficulties left geometers wishing for a simpler, less intricate, and more direct solution. The distinguished \textsc{Legendre} did not fully satisfy this wish with his new proof of the principal theorem (\textit{Hist. de l’acad. roi. des sc.} 1788, \textit{Sur les intégrales doubles}), although his exquisite analytical skill deservedly won the approval of all geometers\HistoricalNote{008}. Later the eminent \textsc{Biot} tried to simplify one solution, and the eminent \textsc{Plana} the other (\textit{Mém. de l’institut T.} VI, \textit{Memorie di matematica e di fisica della società italiana T.} XV). Even so, anyone will readily grant that both solutions rank among the most intricate applications of analysis.

We therefore hope that analysts and astronomers will welcome a new solution of this much-celebrated problem that proceeds along an entirely different path and, if we are not mistaken, has a simplicity that leaves nothing further to be desired.

Our solution itself will take up very few pages. Before turning to the problem itself, to which this paper is devoted, we nevertheless think it worthwhile to carry out some preliminary investigations somewhat more generally and explain them at greater length than our purpose in itself would require, since they can also be applied to advantage on other occasions.
